\section{Summary of Contributions}
This thesis studied whether a multimodal stress detector can avoid relying on a spurious facial attribute, in two lines of work. Its main results are as follows.
\begin{enumerate}
    \item \textbf{Physiological information in facial video (Line A).} A pre-registered test on four UBFC-Phys participants found recoverable pulse information in facial video (POS, $p = 0.0314$), with consistent exploratory results (CHROM, $p = 0.0001$; PBV, $p = 0.0103$). A visual branch trained on real video therefore has access to genuine as well as spurious cues.
    \item \textbf{Orthogonal subspace decomposition.} EQUITAS-RCMF projects visual features onto orthogonal causal and confounder subspaces by a thin QR projection applied at every forward pass, with an orthogonality error of $1.7149 \times 10^{-6}$.
    \item \textbf{Invariance to the scar (Line B).} On a benchmark in which the scar accompanies stress in 85\% of training samples, the deployed model's accuracy was 0.617 whether the scar was positively associated with stress, independent of it, or negatively associated with it, with a counterfactual gap of 0.00064 and an equalized-odds gap of at most 0.065.
    \item \textbf{Penalty alone is insufficient.} A model trained with a counterfactual penalty but without the subspace decomposition reduced the counterfactual gap to 0.0019 but reached only 50.20\% accuracy, close to chance.
    \item \textbf{Deployment.} The compiled model does not require the scar mask and runs at 0.90~ms per frame (1112.7 frames per second) on a desktop CPU.
\end{enumerate}

\section{Boundaries of the Results}
\subsection{Cohort Size}
Both lines rest on small cohorts. Line~A uses four participants, and in Line~B every test label comes from three WESAD participants. A network with over a million parameters, trained on 2,344 windows from ten people, can easily memorize a spurious feature, which makes the setting a demanding test of invariance, but it also limits what can be concluded about variation across people. Evaluation on the full UBFC-Phys cohort of 56 participants, and with cross-validation over all 15 WESAD participants, is needed before the results can be generalized.

\subsection{Invariance and Accuracy}
In the Line~B benchmark, the face images carry no stress information, so a model that ignores the visual input entirely is behaving correctly, and the best achievable accuracy is that of a model using physiology alone. The central open question is therefore whether EQUITAS-RCMF's accuracy of 0.617 matches that of a physiology-only model on the same split, in which case its invariance comes at no cost, or falls below it, in which case the fairness mechanisms reduce accuracy. This comparison will be reported in the final version.

\subsection{Focus-Scale Difference}
The gate is trained with the privileged focus, which is unbounded, but deployed with the autonomous estimate, which lies between 0 and 1. The deployed gate may therefore penalize scar-focused visual features less strongly than the trained gate. The reported autonomous-mode results reflect the deployed behaviour, but matching the two signals, for example by training the autonomous estimate to reproduce the privileged focus, is a necessary refinement.

\subsection{Quantization and Orthogonality}
Standard post-training quantization rounds weights to a low-precision grid, which would perturb the projection matrix so that its rows are no longer exactly orthonormal. The effect of this perturbation on the separation between the causal and confounder subspaces has not yet been measured. Quantization methods that preserve orthogonality, or quantization-aware training with an orthogonality constraint, are therefore a necessary direction before the model is deployed in reduced precision.

\section{Future Work}
\subsection{Evaluation}
The most immediate extensions are to evaluate a physiology-only model and the naive ERM model on the same split and regimes as EQUITAS-RCMF; to report the sanity-check results; to evaluate additional training seeds; and to compare EQUITAS-RCMF with group-robust baselines that use the known scar groups~\cite{sagawa2020}.

\subsection{Combining the Two Lines}
Line~A shows that real facial video carries physiological information, whereas the Line~B benchmark deliberately uses face images that carry none. The natural next step is a benchmark built from real video, such as UBFC-Phys, with a synthetic scar added under a controlled association with the label. That setting would test whether EQUITAS-RCMF can ignore the spurious attribute while still using the legitimate visual evidence, which the present benchmark cannot test.

\subsection{Alternative Backbones}
The subspace decomposition does not depend on a particular backbone. Vision Mamba replaces self-attention with a bidirectional state-space model whose cost grows linearly rather than quadratically with the number of image patches~\cite{zhu2024vim}, which could make higher-resolution facial crops feasible on constrained hardware, although efficient implementations of its scan operation currently target GPUs, and its advantage on mobile CPUs has not been established. RepViT incorporates design elements of vision transformers into a mobile convolutional network through structural re-parameterization, with a reported ImageNet top-1 accuracy above 80\% at about 1.0~ms latency on an iPhone~12~\cite{wang2024repvit}. Replacing the MobileNetV3-Small backbone with either architecture would test whether the approach generalizes across backbone designs.

\section{Concluding Remarks}
This thesis set out to design an edge-deployable model whose stress predictions do not depend on a known spurious visual attribute. On a controlled benchmark, EQUITAS-RCMF met that aim: its accuracy and predictions were unchanged when the association between the scar and the label was reversed, it ran well within real-time limits, and it did not require the scar annotation at deployment. The thesis also showed, in real video, that the visual modality carries genuine physiological information. Two questions remain open: whether the model's invariance comes at a cost in accuracy relative to physiology alone, and whether it can be achieved while also using legitimate visual evidence. Answering them, and extending the evaluation to larger cohorts, will determine whether structural constraints of this kind can make multimodal biometric systems both fair and useful.
